% Generated by Listas/make_statements.py from pratica-lecture-06.md, item 20. Do not edit.
\textbf{Source.} Muñoz-Garcia \& Toro-Gonzalez (2019), \emph{Strategy and Game Theory: Practice Exercises with Answers}, Ch. 6, Exercise 5 ``Collusion When N Firms Compete in Prices'', pp.~290--291 {[}PDF 301--302{]}. Reproduced verbatim (statement only; the book's worked answer follows it).



\begin{quote}

Consider a homogenous industry where \(n\) firms produce at zero cost and play the Bertrand game of price competition for an infinite number of periods. Assume that:



\begin{itemize}

\tightlist

\item

  When firms choose the same price, they earn a per-period profit \(\pi(p) = p\alpha\frac{D(p)}{n}\), where parameter \(\alpha\) represents the state of demand.

\item

  When a firm \(i\) charges a price of \(p_i\) lower than the price of all of the other firms, it earns a profit \(\pi(p_i) = p_i\alpha D(p_i)\), and all of the other firms obtain zero profits.

\end{itemize}



Imagine that in the current period demand is characterized by \(\alpha = 1\), but starting from the following period demand will be characterized by \(\alpha = \theta\) in each of the following periods. All the players know exactly the evolution of the demand state at the beginning of the game, and firms have the same common discount factor, \(\delta\).



Assume that \(\theta > 1\) and consider the following trigger strategies. Each firm plays the monopoly price \(p_m\) in the first period of the game, and continues to charge such a price until a profit equal to zero is observed. When this occurs, each firm charges a price equal to zero forever. Under which conditions does this trigger strategy represent a SPNE? {[}Hint: In particular, show how \(\theta\) and \(n\) affects such a condition, and give an economic intuition for this result.{]}

\end{quote}

